\chapter{Tutorial 3: hot bands and emission at a given temperature}\label{ch:tut-thermal}

% TODO: sampling of the starting states (Monte Carlo / thermal), weights
% and nrepeats, energy bins (Ehist*), anharmonic energy histograms;
% splitting the work into parts and jobs (doc/polyad-jobs.md, T1/T2);
% example on pyrene at a moderate temperature with a small sample.
This chapter will show how to compute the spectrum of molecules with a
given internal energy or temperature, by averaging over a sample of
starting states, and how to split the work into cluster jobs.

The temperature enters through \opt{kT}, in cm$^{-1}$. It is convenient to
give the temperature in kelvin as a variable (Section~\ref{sec:inp}):
\begin{lstlisting}
temperature = 500        # K
kT = 0.6950348*temperature
\end{lstlisting}
